%% background.tex — Section 2: Background

\subsection{Institutional Background}
\label{subsec:institutional}

Ukraine is an upper-middle-income country in which primary education spans four years
(grades~1--4), followed by basic secondary (grades~5--9) and upper secondary schooling.
The Ministry of Education and Science sets national curricula and standards, while
administration and financing are shared with sub-national governments whose fiscal
role expanded following decentralisation reforms in the 2010s. Public education
expenditure averaged 4--6~percent of GDP in the years preceding the 2022 invasion
\citep{unesco2022ukraine}.

Educational access was nearly universal before the pandemic, with primary enrolment
rates exceeding 95~percent. Learning quality, however, was a persistent concern.
In Ukraine's first PISA cycle in 2018, 15-year-old students scored below the OECD
average in reading, mathematics, and science---by margins of 27, 46, and 38~points,
respectively---placing Ukraine in the bottom third of participating countries.
Achievement gaps between urban and rural schools and across socio-economic groups
were substantial, reflecting longstanding disparities in school resources and
household learning environments.

\subsection{COVID-19 Disruption}
\label{subsec:covid}

All Ukrainian schools closed on 12~March 2020 and remained fully remote until the
end of the academic year on 29~May~2020---a nationwide shutdown of 79~days. Schools
reopened in September~2020 under an adaptive-quarantine framework that linked
instructional modality to local epidemiological risk: districts classified green,
yellow, or orange could operate in person, but automatically reverted to remote
instruction upon entering the red zone. Because classifications were revised
frequently, realised exposure to distance learning varied substantially across
regions and, within the same region, across schools.

UNESCO's global school-closure tracker records approximately 18~weeks of full closure
and 13~weeks of partial or hybrid operation for Ukraine between February~2020 and
March~2022---roughly 41~weeks, or about 1.1 school years, of disrupted instruction
\citep{unesco2025closures}. These figures reflect official policies; they overstate
disruption for schools that maintained in-person instruction under the adaptive
regime and understate it for those that exceeded policy thresholds in practice.

Table~\ref{tab:timeline} summarises the key dates: school-year starts,
monitoring assessments, and COVID-19 disruptions relative to the two waves
used in this study.

\begin{table}[htbp]
\centering
\caption{Timeline of school years, monitoring assessments, and COVID-19 disruptions}
\label{tab:timeline}
\begin{tabular}{ll}
\toprule
\textbf{Date} & \textbf{Event} \\
\midrule
1 Sep 2017          & Start of 2017/18 school year (cohort assessed in 2018) \\
17 Apr--18 May 2018 & First monitoring assessment (366 schools, 486 classes) \\
\addlinespace
1 Sep 2019          & Start of 2019/20 school year (pandemic strikes mid-year) \\
12 Mar--29 May 2020 & Nationwide full closure (79 days; $\approx$11 weeks) \\
\addlinespace
1 Sep 2020          & Start of 2020/21 school year under adaptive quarantine \\
8--24 Jan 2021      & Winter lockdown (17 days; $\approx$2.5 weeks) \\
25 Jan--15 Apr 2021 & Adaptive quarantine resumes \\
16 Apr--17 May 2021 & Second monitoring assessment (355 schools, 475 classes) \\
\bottomrule
\end{tabular}
\end{table}

During the 2021 monitoring cycle, teachers reported the total duration of remote
instruction their class experienced over the 2020/21 school year. I use these
responses to construct a class-level measure of realised distance-learning exposure.
Table~\ref{tab:remote_dist} reports the distribution for the full 2021 student sample.

\begin{table}[htbp]
\centering
\caption{Distribution of 2021 students by teacher-reported remote learning duration}
\label{tab:remote_dist}
\begin{tabular}{lrr}
\toprule
\textbf{Duration of remote learning} & \textbf{Students} & \textbf{Share (\%)} \\
\midrule
0 months          &   787 &  9.9 \\
Less than 1 month & 2,311 & 29.2 \\
1--2 months       & 3,819 & 48.2 \\
3--4 months       &   818 & 10.3 \\
5 or more months  &   189 &  2.4 \\
\midrule
Total             & 7,924 & 100.0 \\
\bottomrule
\end{tabular}
\begin{tablenotes}
\small
\item \textit{Notes}: Teacher questionnaire responses from the 2021 monitoring wave.
Duration refers to the 2020/21 school year only and excludes the spring~2020 closure.
\end{tablenotes}
\end{table}

Nearly half of students were in classes with 1--2~months of remote instruction during
the 2020/21 year, while 13~percent experienced three months or more. A further
10~percent reported zero remote learning, likely reflecting schools that remained
continuously in the green or yellow zone. This distribution reveals a wide gap
between the aggregate 41-week policy estimate and realised classroom-level exposure,
and motivates the heterogeneity analysis in Section~\ref{sec:heterogeneity}.
